\section{Results}\label{sec:results}

This section presents the experimental results for our agentic fraud detection system. We evaluated performance across multiple dimensions including classification accuracy, robustness, drift detection, and adaptation.

\subsection{Classification Performance}

Our system achieves strong results on the fraud detection task. The model reached an F1 score of 0.9877, indicating high precision and recall balance. The ROC-AUC score of 0.9869 shows excellent ability to distinguish between fraud and normal transactions across different thresholds.

The precision and recall both reached approximately 0.99, meaning the model correctly identifies fraudulent transactions while keeping false alarms low. These results demonstrate that combining CTGAN data augmentation with adversarial training produces an effective fraud detector that performs well on real-world data.

\subsection{Robustness under Adversarial Attacks}

We tested how well the model handles adversarial perturbations. We evaluated performance at different attack strengths, ranging from zero to 0.30 perturbation levels.

The model maintained strong performance even under attack. At the highest perturbation level, the F1 score remained above 0.97, only dropping slightly from the clean F1 of 0.9877. This shows the model is robust against attempts to trick it by modifying transaction features.

The overall robustness score exceeded 98 percent, calculated as the ratio of average F1 under attack to clean F1. This far exceeds typical target thresholds. The attack success rate, measuring performance degradation under attack, was minimal.

These results demonstrate that adversarial training effectively prepares the model to handle evasion attempts. Even when fraudsters try to manipulate their transactions to bypass detection, the model maintains high accuracy.

\subsection{Drift Detection and Adaptation}

We tested the system's ability to detect and adapt to concept drift. We trained a model on clean data, then tested on intentionally drifted data to measure degradation and recovery.

For our main experiment:
\begin{itemize}
  \item \textbf{Initial (clean data)}: F1 = 0.9877, ROC-AUC = 0.9869, Drift = No
  \item \textbf{After Drift (no retraining)}: F1 = 0.9726, ROC-AUC = 0.9271, Drift = \textbf{YES}
  \item \textbf{After Adaptation}: F1 = 0.9929, ROC-AUC = 0.9943, Drift = No
\end{itemize}

The drift was detected with PSI = 8.09 (far exceeding our 0.20 threshold) and KS = 0.50 (exceeding our 0.05 threshold). Under drift, the ROC-AUC dropped from 0.9869 to 0.9271, a 6.0 percent degradation.

When drift was detected, the system automatically triggered retraining. The model retrained on the drifted data successfully recovered. The final model achieved F1 = 0.9929, which is better than the original model, demonstrating effective adaptation.

This proves the practical value of our MAPE-K architecture. The LLM decision node correctly identifies when drift requires intervention, and the system can adapt to changing fraud patterns without manual intervention.

\subsection{Classification Performance}

Our system achieves strong results:

\begin{table}[h]
\centering
\caption{Performance Metrics}
\begin{tabular}{|l|c|}
\hline
Metric & Value \\
\hline
F1-Score & 0.9877 \\
ROC-AUC & 0.9869 \\
Precision & 0.99 \\
Recall & 0.99 \\
Robustness & 98.34\% \\
Drift Detection (PSI) & 8.09 \\
Drift Detection (KS) & 0.50 \\
Pattern Accuracy & 100\% \\
\hline
\end{tabular}
\end{table}

These results demonstrate that our agentic MAPE-K system with LLM decision-making provides effective autonomous fraud detection with robust drift detection and adaptation capabilities.